%% kapa: exam=1 round=27 year=2016 period=2 form=A subject=appraisal_law subject_ko=감정평가관계법규 seq=13 number=13
%% source: Q-Net 감정평가사 기출문제 (cst00309) artlSeq=5169402 file=extracted/1008/5169402/제27회 감정평가사 1차 문제/제27회 감정평가사 2교시문제(A형).pdf
%% conversion: kordoc 4.24.51 -> tools/gen1.py (auto); review=proofread-sonnet; flags=none
\begin{question}{13}
국토의 계획 및 이용에 관한 법령상 토지에의 출입 등에 관한 설명으로 옳은 것은?
\begin{choices}
  \item 도시ㆍ군계획에 관한 기초조사를 위해 타인의 토지에 출입하는 행위로 인하여 손실을 입은 자가 있으면, 그 행위자가 손실을 보상하여야 한다.
  \item 도시․군계획시설사업에 관한 조사를 위하여 필요한 경우 행정청인 도시계획시설사업의 시행자는 허가 없이 타인의 토지에 출입할 수 있다.
  \item 도시․군계획시설사업에 관한 조사를 위하여 타인의 토지에 출입하려는 자는 시ㆍ도지사의 허가를 받아야 하며, 토지의 소유자ㆍ점유자 또는 관리인의 동의를 받아야 한다.
  \item 도시ㆍ군계획시설사업의 시행자는 타인의 토지를 임시통로로 일시사용하는 경우 토지의 소유자ㆍ점유자 또는 관리인의 동의를 받을 필요가 없다.
  \item 일출 전이나 일몰 후에는 그 토지 점유자의 승낙여부와 관계없이 택지나 담장 또는 울타리로 둘러싸인 타인의 토지에 출입할 수 없다.
\end{choices}
\end{question}
\answer{2}
